\BourbakiExercises{习题}

\begin{enumerate}
\def\labelenumi{\arabic{enumi})}
\tightlist
\item
  设 \(\mathrm{K}\) 是一个交换域，\(\mathrm{A}\) 是一个中心单 K-代数，\(\mathrm{L}\) 是 \(\mathrm{A}\) 的一个半单交换子代数，\(\mathrm{L}'\) 是它在 \(\mathrm{A}\) 中的交换子。
\end{enumerate}

\begin{enumerate}
\def\labelenumi{\alph{enumi})}
\item
  假设 A 是一个有限维 K-向量空间 V 的自同态代数。证明我们有 \([L:K][L':K] \geqslant [A:K]\)（参见 VIII, 第 78 页，命题 1）；等号成立当且仅当 L-模 V 是自由的。由此推出不等式严格成立的一些例子。
\item
  在一般情形证明不等式 \([L:K][L':K] \geqslant [A:K]\)。
\end{enumerate}

\begin{enumerate}
\def\labelenumi{\arabic{enumi})}
\setcounter{enumi}{1}
\tightlist
\item
  设 K 是一个交换域，A 是 K-代数 \(\mathbf{M}_{5}(\mathrm{K})\)，B 是由形如
\end{enumerate}

\[
\left( \begin{array}{c c c} x & 0 & 0 \\ 0 & Y & 0 \\ 0 & 0 & Y \end{array} \right)  ,
\]

的矩阵所成的子代数，C 是由形如

\[
\left( \begin{array}{c c c c} x & 0 & 0 & 0 \\ 0 & x & 0 & 0 \\ 0 & 0 & x & 0 \\ 0 & 0 & 0 & Y \end{array} \right)  ,
\]

的矩阵所成的子代数，其中 x 取遍 K 而 Y 取遍代数 \(\mathbf{M}_{2}(\mathrm{K})\)。证明 B 与 C 是包含 A 的中心 Z 的半单代数。然后证明存在一个 K-代数同构 \(f: B \to C\)，它使 Z 的元素不变，但不存在 A 的任何扩张 f 的自同构（注意这样的自同构必须把 B 的一个单分量映到 C 的一个单分量）。

¶ 3) a) 设 K 是一个交换域，A 是一个 K-代数，B 是 A 的一个有限秩的中心单子代数。设 B' 是 B 在 A 中的交换子。证明从 B ⊗K B' 到 A 的典范同态是双射的。（注意 B 与 B' 在 K 上线性分离（VIII, 第 225 页，习题 5, d)）。）进而证明 A 是一族 (B ⊗K B°)-模 (M \(_{t}\) ) \(_{t\in I}\) 的直和，这些模同构于 B，并注意对每个 \(t\in I\)，M \(_{t}\) 中由这样的同构对应于 1 ∈ B 的元素属于 B'。）

\begin{enumerate}
\def\labelenumi{\alph{enumi})}
\setcounter{enumi}{1}
\tightlist
\item
  反之，设 B 是一个 K-代数，使得对每个包含 B 且与 B 有同一单位元的代数 A，从 \(B \otimes_{K} B'\) 到 A 的典范同态是双射的。证明代数 B 是中心单的且是有限秩的（取 \(A = \text{End}_{K}(B)\)；由假设推出 B 是中心的且拟单的，并利用 VIII, 第 225 页的习题 5, e)）。
\end{enumerate}

\begin{enumerate}
\def\labelenumi{\arabic{enumi})}
\setcounter{enumi}{3}
\item
  设 \(\mathrm{K}\) 是一个交换域，\(\sigma\) 是 \(\mathrm{K}\) 的一个异于恒等映射的自同构，\(\mathrm{D}\) 是体 \(\mathrm{K}((\mathrm{X}))_{\sigma}\)（VIII, 第 19 页，习题 21, e)）。设 \(\tau\) 是 \(\mathrm{K}\) 的一个与 \(\sigma\) 交换的自同构。证明存在 \(\mathrm{D}\) 的唯一的自同构，它扩张 \(\tau\) 且固定 X。由此给出一个固定中心且非内的体自同构的例子。
\item
  设 K 是一个交换域，V 是 K 上一个维数为 n 的向量空间，W 是 V 的一个异于 \(\{0\}\) 与 V 的子空间。用 a 表示代数 \(A = \operatorname{End}_{K}(V)\) 中由满足 \(\operatorname{Im} u \subset W \subset \operatorname{Ker} u\) 的自同态 u 组成的子集，用 B 表示 A 的子代数 \(K + a\)。证明 B 是 A 的一个极大交换子代数，它不是半单的，且它在 K 上的次数可以 \textgreater n。
\end{enumerate}

¶ 6) 设 D 是一个体，其中心为 K 且异于 K，使得 D 的每个元素在 K 上是代数的。

\begin{enumerate}
\def\labelenumi{\alph{enumi})}
\item
  证明 D 包含一个次数 \(>1\) 的可分交换扩张。
\item
  设 m 是一个整数；假设 D 的每个元素在 K 上的次数 \(\leqslant m\)。证明我们有 \([D:K]\leqslant m^{2}\)（先用 a) 与 VIII, 第 259 页的定理 5 证明 D 是有限次数的）。
\end{enumerate}

\begin{enumerate}
\def\labelenumi{\arabic{enumi})}
\setcounter{enumi}{6}
\tightlist
\item
  设 A 是一个交换环，M 与 N 是 A-模，\(u: M \to N\) 是一个 A-线性映射。对 A 的每个极大理想 m，用 \(u(\mathfrak{m})\) 表示由 u 诱导的同态 \(M/mM \to N/mN\)。
\end{enumerate}

\begin{enumerate}
\def\labelenumi{\alph{enumi})}
\item
  假设 N 是有限生成的，且对 A 的每个极大理想 m，\(u(\mathfrak{m})\) 是满射的。证明 u 是满射的（对 u 的余核应用 VIII, 第 81 页的引理 1）。
\item
  假设 M 与 N 是有限生成的，N 是投射的，且对 A 的每个极大理想 m，\(u(\mathfrak{m})\) 是双射的。证明 u 是双射的（用同样的方法）。
\item
  假设 M 与 N 是投射的且有限生成的，且对 A 的每个极大理想 m，\(u(\mathfrak{m})\) 是单射的。证明 u 诱导一个从 M 到 N 的一个直因子子模的同构（对同态 \(^{t}u\) 应用 a)）。
\item
  设 B 是一个 A-代数；假设 A-模 B 是忠实的、投射的且有限生成的。证明 A 是 A-模 B 的一个直因子。证明若 B-模同态 \(u_{(B)} : M_{(B)} \to N_{(B)}\) 是双射的，则 u 也是双射的。证明若 B-模 \(M_{(B)}\) 是投射的，则 A-模 M 也是投射的。
\end{enumerate}

\begin{enumerate}
\def\labelenumi{\arabic{enumi})}
\setcounter{enumi}{7}
\tightlist
\item
  设 A 是一个交换环，S 是一个 A-代数；假设 S 是一个投射的、忠实的且有限生成的 A-模。用 E 表示 A-代数 \(S \otimes_{A} S^{o}\)。
\end{enumerate}

\begin{enumerate}
\def\labelenumi{\alph{enumi})}
\tightlist
\item
  证明下列条件等价：
\end{enumerate}

\begin{enumerate}
\def\labelenumi{(\roman{enumi})}
\item
  (E, A)-双模 S 是可逆的。
\item
  A-代数 S 是中心的，且 E-模 S 是投射的。
\item
  典范同态 \(E \longrightarrow \operatorname{End}_{\mathrm{A}}(\mathrm{S})\) 是双射的。
\item
  对 A 的每个极大理想 m，A/m-代数 S/mS 是中心单的。（由 VIII, 第 101 页的定理 1 与 VIII, 第 100 页的推论 2 推出 (i)、(ii) 与 (iii) 的等价性，并由 VIII, 第 252 页的定理 1 与习题 7 推出 (iii) 与 (iv) 的等价性。）
\end{enumerate}

当这些条件满足时，我们说 \(S\) 是 \(A\) 上的一个 Azumaya 代数。

\begin{enumerate}
\def\labelenumi{\alph{enumi})}
\setcounter{enumi}{1}
\item
  若 S 与 T 是 A 上的 Azumaya 代数，则 A-代数 \(S \otimes_{A} T\) 也是。
\item
  设 B 是一个交换 A-代数；若 S 是 A 上的一个 Azumaya 代数，则 B-代数 \(S_{(B)}\) 是一个 Azumaya 代数。若 B-代数 \(S_{(B)}\) 是一个 Azumaya 代数且 B 是一个忠实的且投射的（\(^{*}\)或忠实平坦的\(_{*}\)）A-代数，则 S 是一个 Azumaya A-代数（应用习题 7, d)）。
\end{enumerate}

\begin{enumerate}
\def\labelenumi{\arabic{enumi})}
\setcounter{enumi}{8}
\tightlist
\item
  设 \(\mathrm{A}\) 是一个交换环，\(\mathrm{S}\) 是一个 Azumaya A-代数（习题 8）。
\end{enumerate}

\begin{enumerate}
\def\labelenumi{\alph{enumi})}
\item
  证明 S 的中心（我们把它与 A 等同）是 A-模 S 的一个直因子（利用习题 7, d)）。
\item
  设 M 是一个 (S,S)-双模；用 \(\mathcal{Z}(M)\) 表示 M 的由满足对每个 \(s \in S\) 有 sm = ms 的元素 m 组成的 A-子模。由习题 8, a) 推出典范同态 \(S \otimes_{A} \mathcal{Z}(M) \to M\) 是双射的。
\item
  映射 \(N \mapsto \mathcal{Z}(N)\) 是从 M 的 \((S, S)\)-子双模所成之集到 \(\mathcal{Z}(M)\) 的 A-子模所成之集上的一个递增双射；逆双射把 \(\mathcal{Z}(M)\) 的一个 A-子模 N 映到 M 的 \((S, S)\)-子双模 SN。特别地，映射 \(s \mapsto s \cap A\) 是从 S 的双边理想所成之集到 A 的理想所成之集上的一个递增双射。
\item
  设 \(\rho : \mathrm{S} \to \mathrm{T}\) 是一个 A-代数同态，\(\mathrm{S}'\) 是 \(\rho(\mathrm{S})\) 在 T 中的交换子。证明典范同态 \(\mathrm{S} \otimes_{\mathrm{A}} \mathrm{S}' \to \mathrm{T}\) 是双射的（对 (S,S)-双模 T 应用 \(a\)））。
\item
  假设每个满足对 A 的每个极大理想 m 有 \(\dim_{\mathrm{A/m}}(\mathrm{L}/\mathfrak{m}\mathrm{L}) = 1\) 的投射 A-模 L 都是秩 1 自由的（例如当环 A 是局部环或主理想整环时就是这种情形）。证明 A-代数 S 的每个自同态 f 都是一个内自同构（对 (S,S)-双模 \(S^{f}\) 应用 a)，并利用习题 8, a) 证明 A-模 \(\mathcal{Z}(S^{f})\) 是投射的）。
\end{enumerate}

\begin{enumerate}
\def\labelenumi{\arabic{enumi})}
\setcounter{enumi}{9}
\tightlist
\item
  设 A 是一个交换环，P 是一个忠实的、投射的且有限生成的 A-模。
\end{enumerate}

\begin{enumerate}
\def\labelenumi{\alph{enumi})}
\item
  证明 A-代数 \(\mathrm{S} = \operatorname{End}_{\mathrm{A}}(\mathrm{P})\) 是一个 Azumaya 代数。
\item
  证明左 S-模 P 是投射的。
\item
  设 \(P'\) 是第二个忠实的、投射的且有限生成的 A-模。证明 A-模 \(P \otimes_{A} P'\) 是忠实的、投射的且有限生成的，且代数 \(\text{End}_{\text{A}}(\text{P} \otimes_{\text{A}} \text{P}')\) 可典范地与代数 \(\text{End}_{\text{A}}(\text{P}) \otimes_{\text{A}} \text{End}_{\text{A}}(\text{P}')\) 等同。
\end{enumerate}

¶ 11) 设 A 是一个交换整环，K 是它的分式域，S 是一个 Azumaya A-代数（习题 8）。

\begin{enumerate}
\def\labelenumi{\alph{enumi})}
\item
  证明 K-代数 \(S_{(K)}\) 是中心单的。把 S 与 \(S_{(K)}\) 的一个子代数等同。
\item
  从现在起假设 A 是 K 的唯一一个作为 A-模是有限生成的 A-子代数，\(^{*}\)即环 A 是整闭的，参见 Comm. Alg., V, §1, No.~1, 2*。证明 S 在 \(S_{(K)}\) 的作为 A-模有限生成的 A-子代数中（对包含关系）是极大的（应用习题 9, d)）。
\item
  进而假设 A 是诺特的，且 K-代数 \(S_{(K)}\) 同构于 \(\operatorname{End}_{\mathrm{K}}(\mathrm{V})\)，其中 V 是 K 上一个向量空间。证明存在一个无挠的有限生成 A-模 M，使得 A-代数 S 同构于 \(\operatorname{End}_{\mathrm{A}}(\mathrm{M})\)（设 x 是 V 的一个非零向量；取 M 为当 s 取遍 S 时元素 \(s(x)\) 所成之集）。
\end{enumerate}

\(^{*}d\) 证明在前面的构造中，我们可以假设 M 是自反的（把 M 换成它的双对偶）。进而，若环 A 是正则的（AC, X, §4, n° 2, 第 55 页, définition 1），则 A-模 M 是投射的（应用 AC, X, §4, 第 164 页, exercice 12）。*

*¶ 12) 设 K 是一个交换域，它对一个离散赋值 v（Comm. Alg., VI）是完备的，我们假设该赋值是规范化的（即 \(v(\mathrm{K}^{*}) = \mathbf{Z}\)）。用 A 表示 v 的赋值环，用 \(\kappa\) 表示它的剩余域。回忆（出处同前）若 L 是 K 的一个有限扩张，则赋值 v 在 L 上有唯一的扩张 \(v_{L}\)；我们说该扩张是非分歧的，如果 \(v_{L}\) 取整数值。

\begin{enumerate}
\def\labelenumi{\alph{enumi})}
\item
  设 D 是一个以 K 为中心的体；设 \(Nrd : D \to K\) 是 D 上的约化范数（VIII, 第 340 页，定义 2）。证明映射 \(w_{0} = v \circ Nrd\) 是 D 上的一个赋值（为证明不等式 \(w_{0}(x + y) \geqslant \inf(w_{0}(x), w_{0}(y))\)，化归到 y = 1 的情形，并注意 \(w_{0}\) 在 \(K(x)\) 上的限制是 \(v_{\mathrm{K}(x)}\) 的倍数）。用 w 表示 D 上与 \(w_{0}\) 等价的规范化赋值。
\item
  L 的一个元素 x 在 A 上是整的当且仅当我们有 \(w(x) \geqslant 0\)；这些元素所成之集构成 D 的一个子环 B。
\item
  从现在起假设域 \(\kappa\) 是完满的。证明若 D 异于 K，则它包含 K 的一个异于 K 的非分歧扩张。（在相反的情形，注意对 D 所含的 K 的每个扩张 L 有 \([\kappa_{\mathrm{L}}:\kappa] = 1\)。设 \(b\in \mathrm{B}\)，并设 \(\pi\) 是 D 的一个单值化元。对域 \(\mathrm{K}(b)\) 应用前面的注记，并证明 \(b\) 属于 \(\mathrm{A}(\pi)\) 的闭包。由于子空间 \(\mathrm{K}(\pi)\) 在 D 中是闭的（EVT, I, §2, n° 3, 第 14 页, corollaire 1），我们得到 \(\mathrm{B}\subset \mathrm{K}(\pi)\)，最终得 \(\mathrm{D} = \mathrm{K}(\pi)\)。）d) 证明存在 D 的一个在 K 上非分歧的极大交换子体。（对 D 的约化次数用归纳法，考虑 D 的一个在 K 上非分歧的子域 L，并对 L 在 D 中的交换子应用归纳假设。）*
\end{enumerate}

\begin{enumerate}
\def\labelenumi{\arabic{enumi})}
\setcounter{enumi}{12}
\tightlist
\item
  设 D 是一个体，V 是 D 上一个左向量空间，A 是环 \(\operatorname{End}_{\mathrm{D}}(\mathrm{V})\)，f 是 A 的一个自同构。证明存在 D 的一个自同构 \(\sigma\) 与一个双射映射 \(u: V \to V\)，它关于 \(\sigma\) 是半线性的（II, §1, No.~13, 第 223 页），使得我们有 \(f(a) = u \circ a \circ u^{-1}\)（对每个 \(a \in A\)）。（设 \(V^{f}\) 是以 V 为加法群、以 \((a, v) \mapsto f(a)(v)\) 为外部运算律的 A-模；由
\end{enumerate}

VIII, 第 50 页的命题 4 推出一个 A-同构 \(u: V \to V^{f}\)。然后注意 \(uDu^{-1}\) 等于 D-模 V 的双交换子，即等于 D。）

\begin{enumerate}
\def\labelenumi{\arabic{enumi})}
\setcounter{enumi}{13}
\tightlist
\item
  设 D 是一个体，Z 是它的中心，V 是 D 上一个左向量空间，\(\Omega\) 是交换群 V 的自同态环。给 \(\Omega\) 装备其自然的 (D,D)-双模结构。设 \(\Gamma\) 是 D 的自同构群，\(\Delta\) 是内自同构所成的子群（同构于 \(D^{*}/Z^{*}\)）。
\end{enumerate}

设 S 是 \(\Omega\) 的一个由半线性映射（相对于 \(\Gamma\) 的一个元素的）组成的子集，F 是 \(\Omega\) 的由 S 生成的左 D-线性子空间；它与由 S 生成的右 D-线性子空间重合。

\begin{enumerate}
\def\labelenumi{\alph{enumi})}
\item
  对 \(\theta \in \Gamma / \Delta\)，用 \(S_{\theta}\) 表示 V 到自身的相对于属于 \(\theta\) 的 D 的自同构的半线性映射所成之集，用 \(F_{\theta}\) 表示 F 的由 \(S_{\theta}\) 生成的 D-线性空间。证明 F 是诸子空间 \(F_{\theta}\)（\(\theta \in \Gamma / \Delta\)）的直和。b) 证明若 \(S_{\theta}\) 在 Z 上是自由的，则它在 D 上也是（左和右）自由的（按 V, §6, No.~1, 第 27 页的定理 1 的证明方式推理）。
\item
  设 E 是 F 的一个 (D,D)-子双模。证明 E（在 D 上）由属于 \(F_{\theta}\) 的半线性映射生成。（设 B 是 F 的包含于 S 的一个基。证明 E 由具有下列性质的向量 \(u \in E\) 生成：若 \(u = \sum_{i}^{n} \lambda_{i} b_{i}\)，其中对每个 i 有 \(\lambda_{i} \neq 0\) 与 \(b_{i} \in B\)，则 E 与由诸 \(b_{i}\) 生成的 F 的子空间的交化归为 Du。然后由 a) 推出 u 是半线性的。）
\end{enumerate}

¶ 15) 保留习题 14 的记号；用 A 表示环 \(\mathrm{End}_{\mathrm{D}}(\mathrm{V})\)。设 G 是 A 的一个自同构群，\(\mathrm{G}_0\) 是 G 中是 A 的内自同构的元素所成的子群。用 U(G) 表示 \(\Omega\) 的由 V 的半线性双射 \(s\)（其自同构 \(a \mapsto sas^{-1}\) 属于 G）所生成的（左和右）D-线性子空间，用 \(\mathrm{U}_0(\mathrm{G})\) 表示 A 的由 V 的自同构 \(u\)（其自同构 \(a \mapsto uau^{-1}\) 属于 \(\mathrm{G}_0\)）所生成的 Z-线性子空间。它们都是 \(\Omega\) 的 Z-子代数。

\begin{enumerate}
\def\labelenumi{\alph{enumi})}
\item
  设 \((s_{\alpha})\) 是 \(U_{0}(G)\) 在 Z 上的一个由 \(S_{0}(G)\) 的元素组成的基，并设 \((g_{\beta})_{\beta\in G/G_{0}}\) 是 G 中诸类（mod \(G_{0}\)）的一个代表系。对每个 \(\beta\in G/G_{0}\)，选取一个元素 \(u_{\beta}\)（\(S(G)\) 的），它对应于 \(g_{\beta}\)。证明诸元素 \(s_{\alpha}u_{\beta}\) 构成 \(U(G)\) 在 D 上的一个（左和右）基（利用习题 14, a)）。次数 \([U(G):D]\) 是有限的当且仅当 \([U_{0}(G):Z]\) 与 \((G:G_{0})\) 是有限的；若是这样，则我们有 \([U(G):D]=[U_{0}(G):Z](G:G_{0})\)。
\item
  证明 \(\mathrm{U}_{0}(\mathrm{G})\) 等于 \(\mathrm{U}(\mathrm{G})\cap\mathrm{A}\)，且 \(\mathrm{U}(\mathrm{G}_{0})=\mathrm{D}u_{0}(\mathrm{G})\) 典范地同构于 \(\mathrm{U}_{0}(\mathrm{G})\otimes_{\mathrm{Z}}\mathrm{D}\)（利用习题 14 以及 VIII, 第 227 页的习题 14）。c) \(\mathrm{U}(\mathrm{G})\) 的一个元素是 V 到自身的一个半线性映射当且仅当它形如 \(\lambda su_{\beta}\)，其中 \(\lambda\in D\)、\(s\in\mathrm{U}_{0}(\mathrm{G})\) 且 \(\beta\in\mathrm{G}/\mathrm{G}_{0}\)（利用 a)、b) 与习题 14, a)）。
\item
  把 U(G) 在 \(\Omega\) 中的交换子记作 \(A^{G}\)；它是 A 的由在 G 下不变的元素组成的子环。设 \(Z^{G} = A^{G} \cap Z\)；它是 Z 的在 G 下不变的元素所成的子体。证明 Z 与 \(A^{G}\) 在 \(Z^{G}\) 上线性分离（验证 \(A^{G}\) 在 \(Z^{G}\) 上的一个基是 A 在 Z 上的一个自由子集）。
\item
  证明若 \(U_{0}(G)\) 是一个单环且 Z-代数 \(U_{0}(G)\) 与 D 之一是有限次数的，则环 \(U(G)\) 是拟单的（由 c) 推出 \(U(G_{0})\) 是一个单环，并利用 c) 对 \(U(G)\) 的一个双边理想应用习题 14, c)）。
\item
  假设 V 与 U(G) 在 D 上是有限维的，且环 \(U_{0}(G)\) 是单的。证明环 A{[}G{]} 是单的，且我们有 \([A : A^{G}]_{s} = [U_{0}(G) : Z](G : G_{0})\)（应用 VIII, 第 205 页的命题 12）。证明 \(A^{G}\) 在 \(\Omega\) 中（相应地，在 A 中）的交换子是 U(G)（相应地，\(U_{0}(G)\)）（利用 VIII, 第 139 页的命题 5）。
\end{enumerate}

\(g\) ) 由此推出 A 的固定 \(\mathrm{A}^{\mathrm{G}}\) 的自同构所成的群由诸自同构 \(a\mapsto vav^{-1}\) 组成，其中 \(v\) 取遍 \(\mathrm{U(G)}\) 中的半线性双射所成之集。证明它由 G 与 \(\mathrm{U_0(G)}\) 的可逆元所相伴的内自同构生成（利用 \(c\)）。

\begin{enumerate}
\def\labelenumi{\alph{enumi})}
\setcounter{enumi}{7}
\tightlist
\item
  若 \(G = G_{0}\)，则环 \(A^{G}\) 包含 Z；当 D 在 Z 上是有限维时证明其逆（利用 VIII, 第 257 页的定理 4）。
\end{enumerate}

¶ 16) 保留习题 15 的记号，并进而假设 V 在 D 上是有限维的。称 A 的一个子环 B 在 A 中是弱伽罗瓦的（相应地，伽罗瓦的），如果它在 Ω 中的交换子 B' 作为（左或右）D-向量空间由 V 到自身的半线性映射（相应地，双射）生成。用 DB 表示 Ω 的由 D 与 B 生成的子环。

\begin{enumerate}
\def\labelenumi{\alph{enumi})}
\item
  证明若 B 是单的且弱伽罗瓦的，则 V 是 DB 上的一个有限长度的半单模。（设 W 是 B-模 V 的一个同型分量，S 是 W 的一个单 DB-子模。设 \(\Sigma\) 是属于 \(B'\) 的半线性映射 v 所成之集。注意 \(v(\mathrm{S})\) 是一个 DB-模（对 \(v \in \Sigma\)），并由 VIII, 第 56 页的推论 3 推出 W 包含于诸 \(v(\mathrm{S})\)（\(v \in \Sigma\)）之和。）
\item
  假设环 B 是单的且弱伽罗瓦的，且它在 A 中的交换子 \(B_{0}^{\prime}\) 是单的。证明 V 是一个同型 DB-模（注意 \(B_{0}^{\prime}\) 是 DB 在 \(\Omega\) 中的交换子，并利用 VIII, 第 85 页的命题 5）。由此推出 B 是伽罗瓦的。（设 v 是 \(\Sigma\) 的一个非零元素，且 \(V = \bigoplus_{i=1}^{n} V_{i}\) 是 V 作为同构的单 DB-模的直和的一个分解，使得 \(W_{1} = v(V_{1})\) 非零。注意存在单 DB-模 \(W_{i}\)（对 \(2 \leqslant i \leqslant n\)），同构于 \(W_{1}\)，使得 \(V = \bigoplus_{i=1}^{n} W_{i}\)，并由此推出一个半线性双射 \(w \in B^{\prime}\)，它在 \(V_{1}\) 上与 v 重合且使 \(vw^{-1}\) 属于 \(B_{0}^{\prime}\)。注意单环的每个元素都是可逆元之和，由此得结论。）
\end{enumerate}

¶ 17) 保留习题 15 的记号，并假设 V 在 D 上是有限维的。称 A 的一个自同构群 G 是伽罗瓦的，如果它满足下列条件：

\begin{enumerate}
\def\labelenumi{(\roman{enumi})}
\item
  \((\mathrm{G}:\mathrm{G}_{0})\) 与 \([\mathrm{U}_{0}(\mathrm{G}):\mathrm{Z}]\) 是有限的。
\item
  环 \(U_{0}(G)\) 是单的。
\item
  对 \(\mathrm{U}_{0}(\mathrm{G})\) 的每个可逆元 s，自同构 \(a \mapsto sas^{-1}\) 属于 \(G_{0}\)。
\end{enumerate}

设 \(G\) 是这样一个群。

\begin{enumerate}
\def\labelenumi{\alph{enumi})}
\item
  证明映射 \(H \mapsto A^{H}\) 是从 G 的伽罗瓦子群所成之集到 A 的包含 \(A^{G}\) 且其在 A 中的交换子 \(B_{0}^{\prime}\) 是单的单子环所成之集上的一个递降双射。（由习题 15, f) 与 g) 推出该映射是完全定义的且是单射的。设 B 是 A 的一个满足上述条件的单子环，H 是 A 的固定 B 的自同构群。利用习题 14, c) 与习题 16 证明 B 是伽罗瓦的，再利用 VIII, 第 205 页的命题 12 证明我们有 \(B = A^{H}\) 且 H 是伽罗瓦的。）
\item
  设 \(B_{1}\) 与 \(B_{2}\) 是 A 的包含 \(A^{G}\) 且其在 A 中的交换子是单的单子环，\(\varphi\) 是从 \(B_{1}\) 到 \(B_{2}\) 的一个固定 \(A^{G}\) 的同构。证明存在 G 的一个元素，它与 \(\varphi\) 在 \(B_{1}\) 上重合。
\end{enumerate}

（首先注意我们有 \(\operatorname{long}_{\mathrm{B}_{1}}(\mathrm{V}) = \operatorname{long}_{\mathrm{B}_{2}}(\mathrm{V})\)，这要用到等式 \(h(\mathrm{B}_{1}, \mathrm{A}^{\mathrm{G}}) = h(\mathrm{B}_{2}, \mathrm{A}^{\mathrm{G}})\)。设 W 是 \(\Omega\) 中满足 \(f(bx) = \varphi(b)f(x)\)（对 \(b \in B_{1}\)、\(x \in V\)）的元素 f 所成之集。证明 W 不化归为 0，然后对它应用习题 14, c) 以证明它包含一个非零的半线性映射 w。证明若 S 是 V 的一个不含于 Ker w 的单 \(DB_{1}\)-子模，则 w 诱导一个从 S 到其像的双射，且我们有 \(\operatorname{long}_{\mathrm{B}_{1}}(\mathrm{S}) = \operatorname{long}_{\mathrm{B}_{2}}(w(\mathrm{S}))\)（利用习题 16, b)）。按出处同前的方式推理，得出 W 包含一个在 S 上与 w 重合的半线性双射。）

\begin{enumerate}
\def\labelenumi{\alph{enumi})}
\setcounter{enumi}{2}
\tightlist
\item
  把 a) 与 b) 的结果特殊化到下列情形：
\end{enumerate}

\begin{enumerate}
\def\labelenumi{(\roman{enumi})}
\item
  \(G = G_{0}\)（参见 VIII, 第 259 页，定理 5）。
\item
  \(\mathrm{G}_0 = \{1\}\)。（证明 A 的每个有限自同构群都是伽罗瓦的，且 A 的每个包含 \(\mathrm{A}^{\mathrm{G}}\) 的单子环在 A 中以 Z 为交换子。）
\item
  V 的维数为 1，于是 A 是同构于 \(D^{o}\) 的一个体。（环 \(A^{G}\) 是一个体，且 A 的每个包含 \(A^{G}\) 的子体都以一个体为交换子。研究 D 是交换的这一特殊情形（参见 V, §10, No.~7）。）
\end{enumerate}

\begin{enumerate}
\def\labelenumi{\arabic{enumi})}
\setcounter{enumi}{17}
\tightlist
\item
  设 \(\mathrm{K}\) 是一个交换域，\(\mathbf{Z}\) 是 \(\mathrm{K}\) 的一个次数为 \(n \geqslant 3\) 的循环扩张（V, §11, No.~6, 第 85 页），\(\sigma\) 是 \(\mathbf{Z}\) 在 \(\mathbf{K}\) 上的伽罗瓦群的一个生成元。用 A 表示环 \(\mathbf{M}_2(\mathbf{Z})\)，用 B 表示 A 的由诸矩阵 \(\left( \begin{array}{cc}z & 0\\ 0 & \sigma (z) \end{array} \right)\)（\(z\in \mathbf{Z}\)）组成的子环。a) 保留习题 17 的记号。证明存在 A 的自同构的一个伽罗瓦群 \(\mathbf{G}\)，使得 \(\mathrm{A}^{\mathrm{G}} = \mathrm{K}\)、\((\mathrm{G}:\mathrm{G}_0) = n\) 且 \(\mathrm{U}_0(\mathrm{G}) = \mathrm{A}\)。
\end{enumerate}

\begin{enumerate}
\def\labelenumi{\alph{enumi})}
\setcounter{enumi}{1}
\item
  证明 B 是 A 的一个单子环，但它在 A 中的交换子 \(B_{0}^{\prime}\) 同构于 \(Z \times Z\)。
\item
  设 H 是 A 的固定 B 的自同构群。证明 \(A^{H}\) 等于 \(B_{0}^{\prime}\)。构造一个从 B 到 Z 的同构，它固定 K 但不能扩张为 A 的一个自同构（参见习题 17, b)）。
\item
  证明 A 的子环 B 是弱伽罗瓦的但不是伽罗瓦的（习题 16；描述 B 在 \(\Omega\) 中的交换子）。
\end{enumerate}

¶ 19) 设 D 是一个体，Z 是它的中心，V 是 D 上一个向量空间。把 D 与 V 的位似子体等同。设 A 是 \(\mathrm{End}_{\mathrm{D}}(\mathrm{V})\) 的一个包含 Z 的稠密子环（VIII, 第 129 页，习题 3）。设 E 是 D 的一个包含 Z 的子体，F 是它在 D 中的交换子。证明 Z-代数 \(\mathrm{A} \otimes_{\mathrm{Z}} \mathrm{E}\) 同构于 \(\mathrm{End}_{\mathrm{F}}(\mathrm{V})\) 的一个稠密子代数（利用 VIII, 第 227 页的习题 14，并考虑 AE 在 \(\mathrm{End}_{\mathrm{Z}}(\mathrm{V})\) 中的交换子）。当 E 是 D 的一个极大子体时就是这种情形。

